% ===================================================================
%  BẢNG Số 15 — Chợ. --- Thức ăn.
%
%  Traduction, faite en 2026, en vietnamien standard d'aujourd'hui,
%  sur l'ido de texto/io. Regles de la colonne : voir 10-bang-01.tex.
%
%  LE MARCHE EST LE TABLEAU LE PLUS VIETNAMIEN DE TOUT LE LIVRET,
%  DEVANT MEME LA COUR DE FERME. Chợ le marche, hàng l'etal et celui
%  qui le tient, cá le poisson, thịt la viande, rau le legume, gạo le
%  riz, muối le sel, đường le sucre, dầu l'huile, giấm le vinaigre :
%  cent cinquante mots, et pas dix d'empruntes. La raison est simple
%  et elle vaut d'etre dite : le marche est l'institution que ce pays
%  n'a jamais eu besoin d'importer.
%
%  MAIS IL Y MANQUE DES CHOSES, ET ON NE LES INVENTE PAS. Le
%  camembert, le gruyere, le hareng fume, le cornichon, la truffe et
%  la langouste n'existaient pas ici en 1926 : on garde leur nom
%  francais, transcrit — « ca-măng-be », « gờ-ruy-e », « cá trích
%  hun khói », « dưa chuột muối » — et on le laisse voir. Une
%  traduction qui donnerait a chacun un nom local ferait croire que
%  la chose y etait.
%
%  « ENCAN » ETAIT LE MOT DU QUEBEC POUR LA CRIEE ; ici c'est « bán
%  đấu giá », vendre au prix qui monte. Le compose est transparent et
%  le mot est courant : la vente au plus offrant est une pratique de
%  marche, non une institution etrangere.
%
%  ET LE PIEGE DU QUEBEC — « liqueur », qui y designe une boisson
%  gazeuse — n'en est pas un ici : « rượu mạnh », l'alcool fort, ne
%  se confond avec rien, et l'ido enumere aussitot rhum, cognac,
%  kirsch, anisette et curacao pour lever le moindre doute.
% ===================================================================

%%K t15-apar-1 apar
\VUsaut{4.0mm}
\VUtitre{15.0pt}{60}{BẢNG Số 15}
\VUsaut{3.0mm}

%%K t15-tit-1 sub
\VUpk{11.4pt}{40}{Chợ. --- Thức ăn.}
\VUsaut{3.0mm}

%%K t15-01-1 p
1. --- Phần lớn những người buôn bán cung cấp cho ta thức ăn cần thiết đều tụ họp dưới
\VUgras{nhà lồng} \textsuperscript{(1)} của \VUgras{chợ có mái} hay ở những phố lân cận.

%%K t15-02-1 p
2. --- \VUgras{Hàng thịt lợn} \textsuperscript{(2)}, bán \VUgras{giăm bông}
\textsuperscript{(3)}, \VUgras{dồi tiết} \textsuperscript{(4)}, \VUgras{xúc xích}
\textsuperscript{(5)}, \VUgras{dồi lòng} \textsuperscript{(6)} và \VUgras{thịt ba chỉ
muối} \textsuperscript{(7)}, đứng cạnh \VUgras{bà hàng bơ và phô mai}
\textsuperscript{(8)}, mà quầy bày đầy \VUgras{phô mai Hà Lan} \textsuperscript{(9)} vỏ
đỏ, \VUgras{phô mai ca-măng-be} \textsuperscript{(10)}, \VUgras{bánh phô mai gờ-ruy-e}
\textsuperscript{(11)} và những \VUgras{tảng bơ} \textsuperscript{(12)} tươi.

%%K t15-03-1 p
3. --- Trước mặt bà, một \VUgras{bà hàng rau} \textsuperscript{(13)} mời khách những quả
\VUgras{cà chua} \textsuperscript{(14)} đẹp, những cây \VUgras{súp lơ}
\textsuperscript{(15)}, \VUgras{xà lách} \textsuperscript{(16)}, \VUgras{bắp cải}
\textsuperscript{(17)}, \VUgras{bí ngô} \textsuperscript{(19)}, \VUgras{dưa vàng}
\textsuperscript{(18)} và cả \VUgras{nấm} \textsuperscript{(20)}. Nhưng một \VUgras{người
giao bia}, người đã dừng chiếc \VUgras{xe giao hàng} \textsuperscript{(21)} chở đầy
\VUgras{thùng bia} \textsuperscript{(22)} ngay trước quầy bà, làm khách của bà không lại
gần được. Một chị làm vườn mang đến cho bà một giỏ \VUgras{a-ti-sô}
\textsuperscript{(23)}.

%%K t15-tit-2 sub
\VUsaut{4.0mm}
\VUpk{10.2pt}{40}{Bán và mua.}
\VUsaut{3.0mm}

%%K t15-04-1 p
4. --- Ở chợ, không khí thật náo nhiệt, vì đã đến giờ \VUgras{bán đấu giá}
\textsuperscript{(24)}. Những \VUgras{bà nội trợ} \textsuperscript{(26)} đến đây mua sắm
dừng lại khi đi qua quầy của \VUgras{bà hàng lòng} \textsuperscript{(25)} để mua
\VUgras{lòng bò} hoặc một cái \VUgras{đầu bê}.

%%K t15-05-1 p
5. --- Một \VUgras{người nấu bếp} \textsuperscript{(27)} béo đang mặc cả một con
\VUgras{cá ngừ} rất đẹp ở hàng \VUgras{bà bán cá} \textsuperscript{(28)}, người tùy ý
nâng giá của mình. Bà nhận được rất nhiều \VUgras{lươn}, \VUgras{cá bơn},
\VUgras{cá hồi} và \VUgras{cá chép}. \VUgras{Cá tuyết} và \VUgras{trai} của bà rất tươi.

%%K t15-tit-3 sub
\VUsaut{4.0mm}
\VUpk{10.2pt}{40}{Rẻ và đắt.}
\VUsaut{3.0mm}

%%K t15-06-1 p
6. --- \VUgras{Bà hàng gia cầm} \textsuperscript{(29)}, ngồi cạnh bà, rất có lương tâm.
Khi bà bán một con \VUgras{gà tây}, một con \VUgras{ngỗng}, một con \VUgras{vịt} hay chỉ
một con \VUgras{gà giò}, bà chỉ đòi đúng giá.

%%K t15-07-1 p
7. --- Ở góc quảng trường, trên tầng một, mới đây một \VUgras{người bán vải}
\textsuperscript{(30)} đã mở hàng, nơi khách mua bao giờ cũng chắc tìm được đủ loại
\VUgras{khăn trải bàn}, \VUgras{khăn ăn}, \VUgras{tạp dề} và \VUgras{khăn lau} rất đẹp.
Cùng tầng ấy, một \VUgras{thợ làm dao} \textsuperscript{(31)} đã mở xưởng của mình. Ông
mài \VUgras{dao} cho các bà hàng trong nhà lồng và làm cho những người làm vườn những
cái \VUgras{liềm quắm} bằng thứ \VUgras{thép} cứng nhất.

%%K t15-tit-4 sub
\VUsaut{4.0mm}
\VUpk{10.2pt}{40}{Hàng tạp hóa.}
\VUsaut{3.0mm}

%%K t15-08-1 p
8. --- Dưới xưởng của ông, mới đây đã mở một cửa hàng thực phẩm lớn. Nó không còn là
\VUgras{hàng tạp hóa} \textsuperscript{(32)} nhỏ bé và tầm thường thuở trước, chỉ bán
những thứ thông dụng : \VUgras{chè} \textsuperscript{(33)}, \VUgras{sô-cô-la}
\textsuperscript{(34)}, \VUgras{đường bánh} \textsuperscript{(37)} và \VUgras{xà phòng}
\textsuperscript{(38)}. Ở đó người ta bán đủ thứ : \VUgras{bánh quy khô}, \VUgras{đồ
hộp} \textsuperscript{(35)}, \VUgras{rượu mạnh} \textsuperscript{(36)}, \VUgras{rượu
rum}, \VUgras{rượu cô-nhắc}, \VUgras{rượu anh đào}, \VUgras{rượu hồi} và \VUgras{rượu
cam}. Ở đó cũng tìm được thịt thú rừng : \VUgras{thỏ rừng} \textsuperscript{(39)},
\VUgras{chim hét} \textsuperscript{(40)}, \VUgras{gà gô} \textsuperscript{(41)},
\VUgras{chim cút} \textsuperscript{(42)}, \VUgras{vịt trời} \textsuperscript{(43)} hay
\VUgras{gà lôi} \textsuperscript{(44)}, cũng dễ như tìm những thứ quả xứ nóng như
\VUgras{chuối} \textsuperscript{(46)} và \VUgras{dứa}.

%%K t15-09-1 p
9. --- Một \VUgras{người bán hàng} \textsuperscript{(45)} rất ân cần luôn sẵn sàng đưa
thứ khách muốn : \VUgras{mứt} \textsuperscript{(47)} nho hay mứt mộc qua, \VUgras{mỡ}
\textsuperscript{(48)}, \VUgras{dưa chuột muối} \textsuperscript{(49)} hay \VUgras{cá
mòi} \textsuperscript{(50)} ướp muối.

%%K t15-09-2 p
Điều ấy rất tiện cho những \VUgras{bà khách} \textsuperscript{(51)}, những người tìm
được, bên cạnh các thứ thông dụng nhất, cả những của đầu mùa hiếm nhất và những trái cây
ngon nhất : \VUgras{lê} \textsuperscript{(52)} mọng nước, \VUgras{mận}
\textsuperscript{(53)}, \VUgras{nho} \textsuperscript{(54)}, \VUgras{đào}
\textsuperscript{(55)}, \VUgras{hạnh nhân} \textsuperscript{(56)} và \VUgras{hồ đào}
\textsuperscript{(57)}.

%%K t15-tit-5 sub
\VUsaut{4.0mm}
\VUpk{10.2pt}{40}{Khách hàng.}
\VUsaut{3.0mm}

%%K t15-10-1 p
10. --- \VUgras{Gạo} \textsuperscript{(58)} và \VUgras{đậu lăng} \textsuperscript{(59)}
để cạnh thùng \VUgras{cá trích} \textsuperscript{(60)} hun khói; \VUgras{khoai tây}
\textsuperscript{(61)} và \VUgras{đậu} \textsuperscript{(63)} nằm bên \VUgras{táo}
\textsuperscript{(62)}, \VUgras{sung} \textsuperscript{(64)}, \VUgras{mận khô}
\textsuperscript{(65)} và \VUgras{hạt phỉ} \textsuperscript{(66)}. Trong cửa hàng ấy có
\VUgras{trứng} \textsuperscript{(67)} tươi và \VUgras{ô liu} \textsuperscript{(68)}; vì
thế những chiếc \VUgras{giỏ} \textsuperscript{(70)} và những chiếc \VUgras{túi lưới đi
chợ} \textsuperscript{(73)} đầy lên rất nhanh.

%%K t15-11-1 p
11. --- \VUgras{Cà phê} \textsuperscript{(72)} của hiệu tuyệt hảo ấy đã được các
\VUgras{chị giúp việc} \textsuperscript{(69)} và các \VUgras{chị bếp} khen ngợi xứng
đáng, vì chính \VUgras{ông chủ hiệu tạp hóa} \textsuperscript{(71)} già rang nó rất cẩn
thận.

%%K t15-tit-6 sub
\VUsaut{4.0mm}
\VUpk{10.2pt}{40}{Cảnh phố chợ.}
\VUsaut{3.0mm}

%%K t15-12-1 p
12. --- Không phải ai đến chợ cũng để mua sắm. Trong dòng người nên thơ trên con ngõ dẫn
tới nhà lồng, ta thấy đủ hạng người. Bên cạnh một \VUgras{người làm hàng thịt}
\textsuperscript{(75)} vui vẻ đang đẩy trước mặt chiếc \VUgras{xe cút kít}
\textsuperscript{(74)}, kìa hai người lính nghỉ phép rất tự hào khoe bộ quân phục sáng
chói của mình : anh \VUgras{kỵ binh nhẹ} \textsuperscript{(76)} với chiếc \VUgras{áo
dôn-man} xanh và chiếc \VUgras{mũ trụ} có chùm lông, và anh \VUgras{cai lính Bắc Phi},
với chiếc quần phồng và cái đầu đội \VUgras{mũ phớt đỏ}.

%%K t15-13-1 p
13. --- \VUgras{Người làm bánh mì} \textsuperscript{(80)}, đứng trên bậc cửa \VUgras{lò
bánh mì} \textsuperscript{(78)}, vừa nhào xong bột \VUgras{bánh mì}
\textsuperscript{(79)} của mình và lấy ra khỏi lò những chiếc \VUgras{bánh sừng bò}
\textsuperscript{(81)}, những chiếc \VUgras{bánh mì sữa} \textsuperscript{(82)} và những
chiếc \VUgras{bánh mì tròn} \textsuperscript{(83)} đang bày trong tủ kính.

%%K t15-14-1 p
14. --- Trên lò bánh mì ở một \VUgras{cô thợ may đồ vải} \textsuperscript{(84)} khéo tay,
người thuê nhiều \VUgras{cô thợ thêu} \textsuperscript{(85)}. Cạnh cô ở một \VUgras{chị
thợ là} \textsuperscript{(86)}, người hồ \VUgras{đồ vải} \textsuperscript{(88)} sau khi
giặt và phơi khô, rồi là nó bằng chiếc \VUgras{bàn là} \textsuperscript{(87)} của mình.

%%K t15-tit-7 sub
\VUsaut{4.0mm}
\VUpk{10.2pt}{40}{Thịt, cá và rau.}
\VUsaut{3.0mm}

%%K t15-15-1 p
15. --- Ở \VUgras{hàng thịt} \textsuperscript{(89)}, đặt ở tầng trệt, người ta bán đủ
loại \VUgras{thịt} : \VUgras{thịt cừu} \textsuperscript{(90)}, \VUgras{thịt bò}
\textsuperscript{(94)} hay \VUgras{thịt bê} \textsuperscript{(92)}, dưới dạng \VUgras{đùi
cừu}, \VUgras{bít tết}, \VUgras{thịt quay} và \VUgras{sườn}. \VUgras{Người làm hàng thịt}
\textsuperscript{(93)} lọc tất cả những thứ thịt ấy và để riêng ra những \VUgras{xương
có tủy} \textsuperscript{(96)}. Ở cửa hàng thịt có một \VUgras{bà hàng hàu}
\textsuperscript{(97)} bán \VUgras{hàu} \textsuperscript{(98)}. Bà tách vỏ nếu khách muốn.

%%K t15-16-1 p
16. --- Tất cả những người buôn bán ấy đều đóng môn bài hay thuế chỗ; vì thế
\VUgras{người cảnh sát} \textsuperscript{(100)} không thương xót đuổi những \VUgras{bà
bán rau rong} \textsuperscript{(99)} đáng thương, những người cạnh tranh với họ bằng
cách rong ruổi trên những chiếc xe đẩy chở \VUgras{cà rốt}, \VUgras{củ cải đỏ},
\VUgras{củ cải}, \VUgras{tỏi tây} hay những \VUgras{bó măng tây}, \VUgras{đậu Hà Lan}
tươi, \VUgras{rau bina}, \VUgras{rau chua}, \VUgras{cải xoong} và \VUgras{rau mùi tây}.

%%K t15-tit-8 sub
\VUsaut{4.0mm}
\VUpk{10.2pt}{40}{Coi chừng cảnh sát!}
\VUsaut{3.0mm}

%%K t15-17-1 p
17. --- Cậu bé bán \VUgras{tỏi} \textsuperscript{(101)} và \VUgras{hành} biết rất rõ
rằng cậu cũng không được phép bán hàng gần chợ. Cậu bỏ chạy, liều làm đổ giỏ
\VUgras{cua}, \VUgras{tôm càng} và \VUgras{tôm} của \VUgras{người bán}
\textsuperscript{(102)} \VUgras{tôm hùm gai} và \VUgras{tôm hùm}.

%%K t15-18-1 p
18. --- Ai nấy đều nhốn nháo và ồn ào; nhưng điều đó không ngăn ta nghe rõ tiếng rao của
\VUgras{người thợ kính} \textsuperscript{(103)} rong, cũng như tiếng gọi của
\VUgras{người bán than} \textsuperscript{(104)} vừa rời chiếc xe của mình mấy phút để
giao một \VUgras{bao than} \textsuperscript{(105)}.
\VUsaut{6.0mm}
\VUfilet{22mm}
